
This report evaluates whether student-accessible quantum hardware is currently capable of running Grover’s search algorithm on game-like problems. Grover provides a quadratic speed-up for unstructured search, but its practical value depends on oracle depth and hardware noise. Three cases are studied: Lights Out, small Sudoku, and procedural map generation via Wave Function Collapse (WFC), each with a classical baseline and a Grover-based formulation implemented in Qiskit and, where feasible, executed on IBM hardware. In noise-free simulation, Grover amplification is observed as expected. On hardware, performance degrades strongly with circuit depth: Lights Out ($2\times2$) drops from $\approx96.13\%$ (ideal) to $\approx16.42\%$ (IBM), and WFC ($3\times3$) from $99.6\%$ valid to $\approx1.4\%$ valid. Overall, classical structure-aware solvers remain superior at the problem sizes currently runnable. 

